\chapter{Observability}

You cannot operate what you cannot see.
Observability is the difference between guessing and knowing.

\section{Logs}
Logs are for detailed events and debugging.
Use structured logs with stable fields.

\section{Metrics}
Metrics summarize system health.
Track rate, errors, and duration for each endpoint.

\section{Tracing}
Distributed traces show where time is spent across services.
Use tracing to find the true bottleneck.

\section{Dashboards}
Dashboards should answer specific questions:
\begin{itemize}[nosep]
  \item Is the system up?
  \item Is it fast enough?
  \item Are errors increasing?
\end{itemize}

\begin{patternbox}{Pattern: instrument the edge}
Instrument request entry and exit points.
This gives you complete coverage of the user-facing path.
\end{patternbox}

\section{Structured log example}
Structured logs should use stable keys:
\begin{lstlisting}[language=JSON]
{
  "level": "info",
  "request_id": "req_42",
  "user_id": "user_7",
  "endpoint": "POST /projects/:id/tasks",
  "status": 201,
  "latency_ms": 84
}
\end{lstlisting}

\section{RED and USE}
The RED method tracks Rate, Errors, Duration.
The USE method tracks Utilization, Saturation, Errors.
Use RED for request paths and USE for system resources.

\section{SLIs and SLOs}
Define a small set of service level indicators.
Latency p95 and availability are common starting points.
SLOs are targets, not guarantees.

\section{Tracing context}
Tracing requires consistent context propagation.
Include a request id in every log and response.
Ensure background jobs carry the same trace context.

\section{Sampling}
Trace sampling reduces overhead but can hide tail issues.
Use adaptive sampling to capture slow requests.
Keep error traces at a higher sampling rate.

\section{Alert design}
Alerts should map to user impact.
Avoid alerting on internal metrics without context.
If an alert does not trigger an action, remove it.

\section{Dashboard layers}
Create dashboards at three levels:
\begin{itemize}[nosep]
  \item service overview for on-call
  \item endpoint detail for debugging
  \item dependency health for incident response
\end{itemize}

\section{Log retention}
Logs have cost and compliance implications.
Define retention policies and restrict access to sensitive logs.

\section{Incident timeline}
Observability is about time.
When incidents happen, you need to see the sequence of events.
Ensure logs, metrics, and traces share timestamps and identifiers.

\section{Error budgets}
Error budgets are the currency of reliability.
If you burn the budget, pause feature work and focus on stability.
This keeps operational work aligned with product priorities.

\section{Examples of SLIs}
Common SLIs include:
\begin{itemize}[nosep]
  \item availability of \code{GET /tasks}
  \item p95 latency for \code{POST /tasks}
  \item background job success rate
\end{itemize}

\section{Instrumentation coverage}
Coverage should include all critical paths:
\begin{itemize}[nosep]
  \item API request entry and exit
  \item database queries
  \item external API calls
  \item queue publish and consume
\end{itemize}

\section{Metrics naming}
Use consistent metric names and labels.
Avoid high-cardinality labels like user id.
High cardinality makes metrics expensive and hard to query.

\section{Example metrics}
\begin{lstlisting}[language=JSON]
{
  "metric": "http_request_duration_ms",
  "labels": {"route": "/tasks", "method": "GET"},
  "value": 42
}
\end{lstlisting}

\section{Trace-to-log correlation}
Traces are most useful when linked to logs.
Include trace and span ids in log fields.
This allows fast navigation from a trace to the exact log lines.

\section{Golden signals}
Use the four golden signals:
\begin{itemize}[nosep]
  \item latency
  \item traffic
  \item errors
  \item saturation
\end{itemize}

\section{Alert fatigue}
Too many alerts reduce response quality.
Review alerts monthly and delete those that do not trigger action.

\section{Runbooks}
Every alert should link to a runbook.
Runbooks should include diagnosis steps and mitigation actions.
If a runbook is missing, the alert is incomplete.

\section{Log levels}
Define when to use \code{debug}, \code{info}, \code{warn}, and \code{error}.
If everything is an error, nothing is.
Consistent levels help reduce noise and improve incident triage.

\section{Log hygiene}
Logs should be concise and actionable.
Avoid dumping entire request bodies.
Redact secrets and personal data before logging.

\section{Correlation ids}
Use a single correlation id per request across services.
Return it in responses so support can ask for it.
Correlation ids turn vague bug reports into precise searches.

\section{Metrics types}
Use counters for rates, gauges for current values, and histograms for latency.
Histograms are essential for p95 and p99 calculations.
Avoid averaging latency, which hides tail behavior.

\section{Histogram buckets}
Pick buckets that match user expectations.
If user timeouts are at 2s, ensure your histogram has buckets around 500ms, 1s, and 2s.
Bad buckets make SLOs hard to interpret.

\section{Trace span design}
Create spans for key operations: handler, database, external calls.
Do not create a span for every function.
Span design is about clarity, not completeness.

\section{Context propagation}
Propagate trace context through queues and background jobs.
If you drop context, traces become useless at service boundaries.
Ensure context flows through async code paths.

\section{Metrics cardinality}
Cardinality can break your metrics system.
Avoid labels like user id, order id, or request id.
If you must label by tenant, limit the number of tenants in metrics.

\section{Exemplars}
Exemplars link metrics to trace samples.
They help you jump from a p99 spike to a real trace.
Use exemplars on latency and error metrics.

\section{Synthetics}
Synthetic probes catch issues before users do.
Run a small set of critical flows on a schedule.
Treat synthetic failures as early warnings, not as noise.

\section{Real user monitoring}
RUM captures client latency and failures.
Client metrics can reveal DNS and network issues that servers cannot see.
If you have a frontend, RUM is a complement to server metrics.

\section{Queue and worker metrics}
Background jobs need their own visibility.
Track queue depth, job latency, and retries.
Alert on growing backlog rather than individual job failures.

\section{Dependency monitoring}
External dependencies should have their own health metrics.
Track error rates and latency for each dependency separately.
This helps you avoid blaming the wrong subsystem.

\section{Alert severity}
Define severity levels and expected response times.
Severity 1 should map to user impact and page on-call.
Severity 3 can be ticketed for later investigation.

\section{On-call handoff}
Shift changes should include a brief operational summary.
Outstanding alerts and known risks should be documented.
Handoffs reduce repeated work and missed context.

\section{Incident review}
Every major incident should produce a timeline and a corrective action list.
Use observability data to validate the timeline.
Review what was missing in metrics or logs and add it.

\section{Example alert rule}
\begin{lstlisting}[language=YAML]
alert: HighErrorRate
expr: rate(http_requests_total{status=~"5.."}[5m]) > 0.05
for: 10m
labels:
  severity: "sev2"
annotations:
  summary: "High 5xx error rate"
\end{lstlisting}

\section{Observability maturity}
Start with logs and metrics.
Add tracing once you have a stable logging and metrics baseline.
Avoid over-investing in tracing before you know what you need to trace.

\section{Cost management}
Observability is not free.
Set budgets for log volume, metric cardinality, and trace sampling.
Cost controls prevent future rollbacks of visibility.

\section{Time synchronization}
All observability depends on accurate time.
Ensure systems use synchronized clocks.
Without time sync, cross-service timelines are misleading.

\section{Service maps}
Service maps show dependencies and call paths.
They are most useful when tied to real traffic data.
Keep maps updated or they quickly become stale.

\section{Health checks}
Expose simple health checks for readiness and liveness.
Readiness should verify dependencies, liveness should be minimal.
Use health checks in deploy pipelines and orchestration systems.

\section{SLO documentation}
Document SLOs alongside the metrics that power them.
SLO docs should include the query and the owner.
SLOs without owners are just numbers.

\section{Example SLO document}
\begin{lstlisting}[language=YAML]
slo: "task_create_availability"
target: 99.9
window: 30d
metric: "http_requests_total"
filter: "route=/tasks method=POST status=2xx"
owner: "productivity-platform"
\end{lstlisting}

\section{Sampling strategy}
Sampling should reflect business importance.
For critical endpoints, sample at a higher rate.
For low-value endpoints, sample aggressively to reduce cost.

\section{Sampling table}
\begin{tabularx}{\textwidth}{@{}lX@{}}
\toprule
Endpoint class & Sampling rate \\
\midrule
Checkout and billing & 50\% \\
Core CRUD & 10\% \\
Background jobs & 5\% \\
\bottomrule
\end{tabularx}

\section{Event logging}
Events such as user signups and payments deserve dedicated logs.
Event logs are different from debug logs and should be structured.
Treat events as data, not just text.

\section{Data retention tiers}
Not all logs need the same retention.
Keep high-value security logs longer than debug logs.
Tiered retention reduces cost without losing critical evidence.

\section{Privacy and compliance}
Observability data can contain personal information.
Define what must be redacted at the source.
Compliance requirements should drive retention and access controls.

\section{Chaos drills}
Use drills to test observability, not just reliability.
If a failure occurs and you cannot see it, instrumentation is insufficient.
Drills reveal blind spots before real incidents do.

\section{Observability reviews}
Review observability alongside feature reviews.
Ask which metrics, logs, and traces should be added for new features.
This keeps instrumentation aligned with product changes.

\section{Debugging workflow}
Start with dashboards, then traces, then logs.
A consistent workflow reduces time to resolution.
Document the workflow so new on-call engineers can follow it.

\section{Log sampling}
High-volume services may need log sampling.
Sample only low-value logs and keep error logs unsampled.
If you sample, include the sampling rate in log metadata.

\section{Instrumentation backlog}
Track missing instrumentation in a backlog.
Close the loop after incidents by adding the missing signals.
Instrumentation work should be planned, not ad hoc.

\section{Release gates}
Set a release gate that checks key SLOs.
If error budgets are exhausted, delay the release.
Release gates tie observability to operational decisions.

\section{Alert tuning loop}
Alerts should be tuned after incidents.
If an alert fired without action, adjust or remove it.
Alert tuning is continuous, not a one-time task.

\section{Dashboard ownership}
Every dashboard should have an owner.
Without owners, dashboards become stale quickly.
Ownership keeps metrics accurate and relevant.

\section{Burn rate alerts}
Use burn rate alerts for SLOs.
Burn rates detect fast and slow budget consumption.
This prevents surprises at the end of the SLO window.

\section{Log search practices}
Standardize log search queries for common incidents.
Store shared queries in runbooks.
This reduces time to diagnose recurring issues.

\section{Exemplar-driven debugging}
If you use exemplars, document how to jump from a metric spike to a trace.
Exemplar workflows save time during high-pressure incidents.
Teach this workflow in on-call training.

\section{Metric naming conventions}
Use a consistent naming scheme for metrics.
Include the unit in the name, such as \code{\_seconds} or \code{\_bytes}.
Consistent naming reduces confusion across teams.

\section{SLI calculation example}
\begin{lstlisting}[language=YAML]
sli:
  name: "task_create_availability"
  good_events: "http_requests_total{route=/tasks,method=POST,status=2xx}"
  total_events: "http_requests_total{route=/tasks,method=POST}"
\end{lstlisting}

\section{Error budget policy}
Define what happens when the error budget is depleted.
Common actions include slowing releases or focusing on reliability work.
Policy clarity prevents debates during incidents.

\section{Playbook template}
\begin{lstlisting}[language=YAML]
playbook:
  alert: "HighErrorRate"
  first_checks:
    - "check recent deploys"
    - "inspect dependency latency"
  mitigations:
    - "rollback"
    - "disable feature flag"
\end{lstlisting}

\begin{checklistbox}{Checklist: observability}
\begin{itemize}[nosep]
  \item Log with stable fields.
  \item Measure latency and error rates.
  \item Add traces for critical flows.
\end{itemize}
\end{checklistbox}
